\documentclass[../../main.tex]{subfiles}
\begin{document}

\paragraph{【新課程】}
{\bf [解]}
$1$から$9$までの数字が書かれた$9$枚のカードから$4$枚を続けて取り出し，左から並べて$4$桁の数をつくる方法はすべてで${}_9\mathrm{P}_4$通りあり，これらは同様に確からしい．
このうち，$1966$より小さくなる数は，千の位（最高位）が$1$で固定され，次の$2$つの場合に分けられる．
\begin{align}
  \begin{cases}
    1^\circ & 1900\text{未満（百の位が8以下）} \\
    2^\circ & 1900\text{以上}1965\text{以下（百の位が9）}
  \end{cases}
\end{align}

\subparagraph{$1^\circ$の場合}
千の位は$1$（$1$通り）で，百の位は$2$から$8$までの$7$通りである．十の位は残りの$7$枚から$1$枚，一の位は残りの$6$枚から$1$枚を選べばよいので，
\begin{align}
  1 \times 7 \times 7 \times 6 = 294 \text{ 通り} \label{eq:1}
\end{align}
である．

\subparagraph{$2^\circ$の場合}
千の位が$1$，百の位が$9$（$1$通り）である．
\begin{itemize}
  \item $1960$未満の場合: 十の位は$2$から$5$までの$4$通りで，一の位は残り$6$枚から$1$枚選ぶので，$4 \times 6 = 24$通り．
  \item $1960$以上$1965$以下の場合: 十の位は$6$（$1$通り）で，一の位は$2, 3, 4, 5$のいずれか（$4$通り）より，$4$通り．
\end{itemize}
合わせて
\begin{align}
  24 + 4 = 28 \text{ 通り} \label{eq:2}
\end{align}
である．\casesend

\cref{eq:1}，\cref{eq:2}より，$1966$より小さくなる数は全部で
\begin{align}
  294 + 28 = 322 \text{ 通り}
\end{align}
ある．したがって，求める確率$P$は
\begin{align}
  P &= \dfrac{322}{{}_9\mathrm{P}_4} \nonumber\\
  &= \dfrac{322}{9 \times 8 \times 7 \times 6} \nonumber\\
  &= \dfrac{46}{9 \times 8 \times 6} \nonumber\\
  &= \dfrac{23}{216} \tag{答}
\end{align}
である．

\paragraph{【旧課程】}
題意より，$l_1$は$x$軸に平行であるから，
\begin{align}
  l_1: y=k \label{eq:3}
\end{align}
とおける．また，
\begin{align}
  \begin{cases}
    C: x+y^2-5=0 \\
    C_2: x^2-y+1=0
  \end{cases} \label{eq:4}
\end{align}
とする．

\begin{figure}[htb]
  \centering
  \begin{tikzpicture}[scale=0.9, >=stealth, every node/.style={font=\small}]

    % --- 1. 座標軸 ---
    \draw[->, thick] (-3.5, 0) -- (6.5, 0) node[right] {$x$};
    \draw[->, thick] (0, -3.5) -- (0, 8.5) node[above] {$y$};
    \node[below left] at (0,0) {$O$};

    % --- 2. 放物線の描画 ---
    % C: x = 5 - y^2 (頂点 (5,0))
    \draw[thick, blue!80!black, domain=-2.8:2.8, samples=80]
    plot ({5 - (\x)^2}, {\x});
    \node[blue!80!black, left] at (-2.8, 2.7) {$C: x+y^2=5$};

    % C_1: x = 5 - (y-6)^2 (頂点 B(5,6))
    \draw[thick, cyan!80!black, domain=3.2:8.8, samples=80]
    plot ({5 - (\x - 6)^2}, {\x});
    \node[cyan!80!black, left] at (-2.8, 8.7) {$C_1: x+(y-6)^2=5$};

    % C_2: y = x^2 + 1 (頂点 A(0,1))
    \draw[very thick, magenta!90!black, domain=-2.5:2.5, samples=80]
    plot ({\x}, {(\x)^2 + 1});
    \node[magenta!90!black, above left] at (-2.3, 6.3) {$C_2: y=x^2+1$};

    % --- 3. 対称軸・折り返し直線 ---
    % l_1: y = 3
    \draw[thick, dashed, orange!90!black] (-3.5, 3) -- (6.5, 3)
    node[right, font=\footnotesize] {$l_1: y=3$};

    % l_2: y = -x + 6
    \draw[thick, dashed, red!80!black] (-1.8, 7.8) -- (6.5, -0.5)
    node[below right, font=\footnotesize] {$l_2: y=-x+6$};

    % 各放物線の軸
    \draw[gray!60, dotted] (-3.5, 6) -- (5, 6); % C1 の軸 y=6

    % --- 4. 主要な点と折り返しの対応関係 ---
    % 頂点 A(0,1), B(5,6), Cの頂点(5,0)
    \coordinate (A) at (0, 1);
    \coordinate (B) at (5, 6);
    \coordinate (V0) at (5, 0);
    \coordinate (M) at (2.5, 3.5); % AB の中点
    \coordinate (L2_axis_intersect) at (0, 6); % y=6 と x=0 の交点

    % C から C_1 への折り返し補助線（頂点同士）
    \draw[dashed, orange!80!black] (V0) -- (B);
    \fill (V0) circle (1.3pt) node[below left=1pt] {$(5,0)$};

    % C_1 から C_2 への折り返し線分 AB
    \draw[thick, gray!70] (A) -- (B);
    \fill[red!80!black] (M) circle (1.2pt) node[above right=1pt, font=\scriptsize] {$M\left(\frac{5}{2},\,\frac{7}{2}\right)$};

    % 直角記号 (AB と l_2 の直交)
    % AB の傾きは 1, l_2 の傾きは -1 なので直交する
    \draw[thick, red!80!black]
    ($ (M) + (-0.15, -0.15) $) --
    ($ (M) + (-0.15, -0.15) + (0.15, -0.15) $) --
    ($ (M) + (0.15, -0.15) $);

    % 頂点プロットとラベル
    \fill[magenta!90!black] (A) circle (1.5pt) node[above left] {$A(0,1)$};
    \fill[cyan!80!black] (B) circle (1.5pt) node[above right] {$B(5,6)$};
    \fill (L2_axis_intersect) circle (1.2pt) node[above right, font=\scriptsize] {$(0,6)$};

    % --- 5. 移動矢印（折り返しのイメージ） ---
    % C -> C_1 (l_1 による鏡映)
    \draw[->, thick, orange!90!black] (4.2, 0.9) to[bend right=25] (4.2, 5.1);
    \node[orange!90!black, font=\scriptsize, fill=white, inner sep=1pt] at (4.7, 3.0) {$l_1$ で対称};

    % C_1 -> C_2 (l_2 による鏡映)
    \draw[->, thick, red!80!black] (3.8, 6.7) to[bend right=25] (1.8, 4.8);
    \node[red!80!black, font=\scriptsize, fill=white, inner sep=1pt] at (3.2, 5.6) {$l_2$ で対称};

  \end{tikzpicture}
  \caption{放物線 $C$ から $l_1$ による鏡映 $C_1$，および $l_2$ による鏡映 $C_2$ への変換}
  \label{fig:parabola_reflections}
\end{figure}

$C$を$l_1$に関して折り返した曲線を$C_1$とすれば，$C$上の点$(x, y)$は$C_1$上の点$(x, 2k-y)$に移るから，
\begin{align}
  C_1: x+(2k-y)^2-5=0 \Longleftrightarrow x+(y-2k)^2-5=0 \label{eq:5}
\end{align}
である．

次に$l_2$について考えよう．
$C_1$の頂点は$B(5, 2k)$，対称軸は$y=2k$である．
一方，$C_2$の頂点は$A(0, 1)$，対称軸は$x=0$である．
$C_1$を直線$l_2$に関して折り返すと$C_2$に重なるから，$C_1$の対称軸$y=2k$と$C_2$の対称軸$x=0$は$l_2$に関して対称でなければならない．
これら$2$直線の交点は$(0, 2k)$であり，対称軸$l_2$はこの交点を通ってかつ傾きが$\pm 1$である．したがって
\begin{align}
  l_2: y = \pm x + 2k \label{eq:6}
\end{align}
とおける．

また，この折り返しにより$C_1$の頂点$B(5, 2k)$は$C_2$の頂点$A(0, 1)$に移るから，線分$AB$の中点$M\left(\dfrac{5}{2}, \dfrac{2k+1}{2}\right)$が$l_2$上にあり，かつ$\overrightarrow{AB} \perp l_2$である．
$\overrightarrow{AB} =
\begin{pmatrix} -5 \\ 1-2k
\end{pmatrix}$，および$l_2$の方向ベクトル$
\begin{pmatrix} 1 \\ \pm 1
\end{pmatrix}$について，複合同順として
\begin{align}
  &
  \begin{cases}
    \dfrac{2k+1}{2} = \pm \dfrac{5}{2} + 2k \\
    \begin{pmatrix} -5 \\ 1-2k
    \end{pmatrix} \cdot
    \begin{pmatrix} 1 \\ \pm 1
    \end{pmatrix} = 0
  \end{cases} \\
  \iff
  &
  \begin{cases}
    1 - \mp 5 = 2k
    -5 \times 1 + (1-2k) \times (\pm 1) = 0
  \end{cases} \\
  \iff
  & k = 3 \quad (\text{複合負})
\end{align}
である．この時，\cref{eq:3,eq:6}より
\begin{align}
  l_1: y=3 \\
  l_2: y=-x+6
\end{align}
となる．

以下十分条件を求める．
$C_1$上の点$P(x, y)$が$l_2: x+y=6$に関する折り返しによって点$Q(x+2X, y+2Y)$に移るとすると，
\begin{align}
  &
  \begin{cases}
    \text{線分$AB$の中点$M$が$l_2$上} \\
    \overrightarrow{PQ} \perp l_2
  \end{cases} \\
  \iff
  &
  \begin{cases}
    (y+Y) = (-x+X)+6 \\
    \begin{pmatrix} X \\ Y
    \end{pmatrix}\cdot
    \begin{pmatrix} 1 \\ -1
    \end{pmatrix} = 0
  \end{cases} \\
  \iff
  & 2X = 2Y = -x-y+6
\end{align}
である．この時$Q$が$C_2$上にある条件から
\begin{align}
  & (x+2X)^2-(y+2Y) + 1 = 0 \\
  & (6-y)^2 - (6-x) + 1 = 0 \\
  & x + (y-6)^2 - 5 = 0
\end{align}
である．
これは$Q$が$l_2$の折り返しにより$P$に移動することを示しており十分である．

以上より求める直線は
\begin{align}
  \begin{cases}
    l_1: y = 3 \\
    l_2: y = -x+6
  \end{cases} \tag{答}
\end{align}
である．

\newpage
\end{document}
